% GENERATED FILE. Edit canonical lessons, then run npm run generate:books.
% canonical-source-hash: fnv1a64:1ca8e904509fed4f
% canonical-lessons: MR-C228-kautuk-karne, MR-C228-pudhe-jane, MR-C228-mage-jane, MR-C228-bacav-karne, MR-C228-nahise-hone

\chapter{To praise, To go forward, To go back, To defend, To disappear}
\label{ch:mr-to-praise-to-go-forward-to-go-back-to-defend-to-disappear}
\hlchaptermodality{228}

\begin{chapteropening}
\textbf{By the end of this chapter:} \emph{I can say to praise, to go forward, to go back, to defend and to disappear.}

Complete the last lesson of chapter 228: I can say to praise, to go forward, to go back, to defend and to disappear.
\end{chapteropening}

\section[\texorpdfstring{\mr{कौतुक} \mr{करणे} (kautuk karṇe)}{कौतुक करणे (kautuk karṇe)}]{\mr{कौतुक} \mr{करणे} (kautuk karṇe) — to praise, to admire}

\label{lesson:MR-C228-kautuk-karne}



\begin{quote}
Before the new one: say the Marathi for to float, then the Marathi for to chase.
\end{quote}

\subsection*{You'll want to know: \mr{कौतुक} \mr{करणे}}
\textbf{\mr{कौतुक} \mr{करणे}} — \emph{kautuk karṇe} — \textquotedblleft{}to praise, to admire\textquotedblright{}.

\begin{grammarlens}[title={What you've built}]
One more word for this chapter.
\end{grammarlens}

\subsection*{Your turn}
\begin{itemize}
\raggedright
  \item \emph{Say it:} \emph{kautuk karṇe}
  \item \emph{Say it:} \emph{kautuk karṇe}, once more
  \item \emph{Say it:} say \emph{kautuk karṇe}
  \item \emph{Recall:} say the Marathi for to float, then the Marathi for to chase, then say \emph{kautuk karṇe} again
\end{itemize}

\begin{tcolorbox}[breakable,colback=teal!4,colframe=teal!35!black,title={Before you move on}]
What does \mr{कौतुक} \mr{करणे} mean? (\textquotedblleft{}To praise, to admire\textquotedblright{}.) Say it once more.
\end{tcolorbox}
\section[\texorpdfstring{\mr{पुढे} \mr{जाणे} (puḍhe jāṇe)}{पुढे जाणे (puḍhe jāṇe)}]{\mr{पुढे} \mr{जाणे} (puḍhe jāṇe) — to go forward}

\label{lesson:MR-C228-pudhe-jane}



\begin{quote}
Before the new one: say the Marathi for to chase, then the Marathi for to praise.
\end{quote}

\subsection*{You'll want to know: \mr{पुढे} \mr{जाणे}}
\textbf{\mr{पुढे} \mr{जाणे}} — \emph{puḍhe jāṇe} — \textquotedblleft{}to go forward\textquotedblright{}.

\begin{grammarlens}[title={What you've built}]
One more word for this chapter.
\end{grammarlens}

\subsection*{Your turn}
\begin{itemize}
\raggedright
  \item \emph{Say it:} \emph{puḍhe jāṇe}
  \item \emph{Say it:} \emph{puḍhe jāṇe}, once more
  \item \emph{Say it:} say \emph{puḍhe jāṇe}
  \item \emph{Recall:} say the Marathi for to chase, then the Marathi for to praise, then say \emph{puḍhe jāṇe} again
\end{itemize}

\begin{tcolorbox}[breakable,colback=teal!4,colframe=teal!35!black,title={Before you move on}]
What does \mr{पुढे} \mr{जाणे} mean? (\textquotedblleft{}To go forward\textquotedblright{}.) Say it once more.
\end{tcolorbox}
\section[\texorpdfstring{\mr{मागे} \mr{जाणे} (māge jāṇe)}{मागे जाणे (māge jāṇe)}]{\mr{मागे} \mr{जाणे} (māge jāṇe) — to go back}

\label{lesson:MR-C228-mage-jane}



\begin{quote}
Before the new one: say the Marathi for to praise, then the Marathi for to go forward.
\end{quote}

\subsection*{You'll want to know: \mr{मागे} \mr{जाणे}}
\textbf{\mr{मागे} \mr{जाणे}} — \emph{māge jāṇe} — \textquotedblleft{}to go back\textquotedblright{}.

\begin{grammarlens}[title={What you've built}]
One more word for this chapter.
\end{grammarlens}

\subsection*{Your turn}
\begin{itemize}
\raggedright
  \item \emph{Say it:} \emph{māge jāṇe}
  \item \emph{Say it:} \emph{māge jāṇe}, once more
  \item \emph{Say it:} say \emph{māge jāṇe}
  \item \emph{Recall:} say the Marathi for to praise, then the Marathi for to go forward, then say \emph{māge jāṇe} again
\end{itemize}

\begin{tcolorbox}[breakable,colback=teal!4,colframe=teal!35!black,title={Before you move on}]
What does \mr{मागे} \mr{जाणे} mean? (\textquotedblleft{}To go back\textquotedblright{}.) Say it once more.
\end{tcolorbox}
\section[\texorpdfstring{\mr{बचाव} \mr{करणे} (bacāv karṇe)}{बचाव करणे (bacāv karṇe)}]{\mr{बचाव} \mr{करणे} (bacāv karṇe) — to defend}

\label{lesson:MR-C228-bacav-karne}



\begin{quote}
Before the new one: say the Marathi for to go forward, then the Marathi for to go back.
\end{quote}

\subsection*{You'll want to know: \mr{बचाव} \mr{करणे}}
\textbf{\mr{बचाव} \mr{करणे}} — \emph{bacāv karṇe} — \textquotedblleft{}to defend\textquotedblright{}.

\begin{grammarlens}[title={What you've built}]
One more word for this chapter.
\end{grammarlens}

\subsection*{Your turn}
\begin{itemize}
\raggedright
  \item \emph{Say it:} \emph{bacāv karṇe}
  \item \emph{Say it:} \emph{bacāv karṇe}, once more
  \item \emph{Say it:} say \emph{bacāv karṇe}
  \item \emph{Recall:} say the Marathi for to go forward, then the Marathi for to go back, then say \emph{bacāv karṇe} again
\end{itemize}

\begin{tcolorbox}[breakable,colback=teal!4,colframe=teal!35!black,title={Before you move on}]
What does \mr{बचाव} \mr{करणे} mean? (\textquotedblleft{}To defend\textquotedblright{}.) Say it once more.
\end{tcolorbox}
\section[\texorpdfstring{\mr{नाहीसे} \mr{होणे} (nāhīse hoṇe)}{नाहीसे होणे (nāhīse hoṇe)}]{\mr{नाहीसे} \mr{होणे} (nāhīse hoṇe) — to disappear}

\label{lesson:MR-C228-nahise-hone}



\begin{quote}
Before the new one: say the Marathi for to go back, then the Marathi for to defend.
\end{quote}

\subsection*{You'll want to know: \mr{नाहीसे} \mr{होणे}}
\textbf{\mr{नाहीसे} \mr{होणे}} — \emph{nāhīse hoṇe} — \textquotedblleft{}to disappear\textquotedblright{}.

\begin{grammarlens}[title={What you've built}]
One more word for this chapter.
\end{grammarlens}

\subsection*{Your turn}
\begin{itemize}
\raggedright
  \item \emph{Say it:} \emph{nāhīse hoṇe}
  \item \emph{Say it:} \emph{nāhīse hoṇe}, once more
  \item \emph{Say it:} say \emph{nāhīse hoṇe}
  \item \emph{Recall:} say the Marathi for to go back, then the Marathi for to defend, then say \emph{nāhīse hoṇe} again
\end{itemize}

\begin{tcolorbox}[breakable,colback=teal!4,colframe=teal!35!black,title={Before you move on}]
What does \mr{नाहीसे} \mr{होणे} mean? (\textquotedblleft{}To disappear\textquotedblright{}.) Say it once more.
\end{tcolorbox}
